% !TEX root = ../main.tex
\begin{abstract}
% arXiv reproduces this as plain text in the listing: no macros, no \cite.
Language model deployments in regulated environments are commonly validated by
comparing outputs against reference answers. This paper argues that the practice cannot
deliver what sign-off requires, and identifies the property that can.

The first of two connected claims is a redefinition of hallucination, indexed to a
decision context. Existing definitions ask whether an assertion is supported by a
supplied source; this one asks the prior question of whether that source governs the
case, by requiring support to be applicable at a resolved context of jurisdiction, valid
time, material facts, and intended use, and by requiring the assertion to have been
produced because of the evidence rather than merely to agree with it. The consequence is
a failure class---authentic, authoritative, correctly cited evidence that does not govern
the decision---which faithfulness and attribution metrics score as clean by construction,
since the entailment genuinely holds. Unlike the truth-conditional definition, this one
can be evaluated at the point of assertion without an answer key, and so functions as a
gate rather than only as a score.

The second claim is that the same architecture makes accuracy uncertifiable. The system
computes a distribution over token sequences conditioned on token sequences and applies
no canonical form, so a correct answer observed once carries no guarantee that a
materially equivalent request will be answered the same way. Accuracy is a property of a
sample and the sample does not generalise. I prove that accuracy places no bound whatever
on the grounding rate---the two vary over the whole unit square, so that neither identifies the other---so no
amount of reference-based evaluation constrains how often a system asserts beyond its
evidence.

The constructive half derives a measurement protocol, nineteen system invariants, a
reference architecture for governed retrieval, and an account of why the division of
ownership across data, policy, platform, validation, and records functions is itself a
principal source of failure.
\end{abstract}
